\section{Discussion}

\paragraph{Generalizability to Other Hardwares.} We argue that the “vector + matrix” architecture of NPUs possesses a certain degree of universality and observe that the boundary between CPUs and NPUs is gradually blurring. Beyond NPUs, modern CPUs have also begun to incorporate dedicated matrix multiplication units, such as Intel AMX and ARM SME, endowing them with a similar "vector + matrix" architecture. Furthermore, we note that modern AI accelerators generally exhibit a significant disparity between general-purpose computing performance and specialized low-precision matrix multiplication capabilities (e.g. NVIDIA GPUs). Although specific hardware architectures may differ, the core ideas behind our techniques maintain broad applicability.

\paragraph{System Performance and Limitations.} (a) Decoding Performance: The current decoding speed of our system is relatively constrained, primarily due to the overhead of dequantization. However, this does not undermine the effectiveness of test-time scaling. Quantized GEMM based on QNN typically utilizes only the DMA and HMX components without introducing HVX computational overhead. Approaches similar to T-MAC~\cite{tmac} could potentially enable efficient GEMV with fine-grained group quantization on NPUs, thereby accelerating the LLM decoding process. (b) Prefill Performance: There remains room for improvement in the prefill performance of our current system. Offloading more operators to the NPU, reducing memory access and communication overhead through operator fusion, and optimizing tiling and pipelining strategies for matrix multiplication could all contribute to enhanced prefill performance. We leave these optimizations to future work. (c) Model Size Constraints: Our current implementation is limited by the 32-bit address space of a single NPU session on older devices. Employing multiple NPU sessions could help alleviate this issue.

\paragraph{Application Scope of Parallel Test-time Scaling.} Although parallel test-time scaling methods currently dominate mathematical reasoning tasks, evidence from recent studies~\cite{yan2025empowering,tot,paranjape2023boosting,hong2023metagpt,hafner2024reasoning,parallel-scaling-law} indicates their extensibility to broader reasoning and planning domains, highlighting substantial generalizable potential.
